% appendix.tex — Complete numerical tables for the cross-format probing study
\section*{Appendix}
\label{sec:appendix}

\setcounter{table}{0}
\renewcommand{\thetable}{A\arabic{table}}

%% ========================================================================
%% Table A1: Complete CKA Values
%% ========================================================================
\begin{table*}[t]
\centering
\caption{Complete CKA similarity values per model, layer, and format pair. Mean is the arithmetic average of the three format pairs, consistent with main text Table~\ref{tab:cka_mean}.}
\label{tab:cka_complete}
\small
\begin{tabular}{ll rrrr}
\toprule
\textbf{Model} & \textbf{Layer} & \textbf{SVG--TikZ} & \textbf{SVG--Asy} & \textbf{TikZ--Asy} & \textbf{Mean} \\
\midrule
Qwen2.5-Coder & L4  & 0.049 & 0.092 & 0.120 & 0.087 \\
Qwen2.5-Coder & L8  & 0.059 & 0.087 & 0.125 & 0.090 \\
Qwen2.5-Coder & L12 & 0.067 & 0.098 & 0.200 & 0.121 \\
Qwen2.5-Coder & L16 & 0.086 & 0.124 & 0.212 & 0.141 \\
Qwen2.5-Coder & L20 & 0.087 & 0.118 & 0.180 & 0.128 \\
Qwen2.5-Coder & L24 & 0.105 & 0.109 & 0.209 & 0.141 \\
Qwen2.5-Coder & L28 & 0.103 & 0.114 & 0.229 & 0.149 \\
\midrule
VisCoder2      & L4  & 0.050 & 0.092 & 0.122 & 0.088 \\
VisCoder2      & L8  & 0.062 & 0.089 & 0.133 & 0.095 \\
VisCoder2      & L12 & 0.075 & 0.103 & 0.218 & 0.132 \\
VisCoder2      & L16 & 0.089 & 0.124 & 0.216 & 0.143 \\
VisCoder2      & L20 & 0.086 & 0.124 & 0.181 & 0.130 \\
VisCoder2      & L24 & 0.102 & 0.125 & 0.208 & 0.145 \\
VisCoder2      & L28 & 0.093 & 0.116 & 0.203 & 0.137 \\
\midrule
Qwen2.5-Base   & L4  & 0.050 & 0.087 & 0.118 & 0.085 \\
Qwen2.5-Base   & L8  & 0.058 & 0.087 & 0.122 & 0.089 \\
Qwen2.5-Base   & L12 & 0.061 & 0.093 & 0.182 & 0.112 \\
Qwen2.5-Base   & L16 & 0.067 & 0.112 & 0.166 & 0.115 \\
Qwen2.5-Base   & L20 & 0.070 & 0.104 & 0.153 & 0.109 \\
Qwen2.5-Base   & L24 & 0.086 & 0.099 & 0.177 & 0.121 \\
Qwen2.5-Base   & L28 & 0.044 & 0.081 & 0.104 & 0.076 \\
\bottomrule
\end{tabular}
\end{table*}

%% ========================================================================
%% Table A2: Procrustes Distance
%% ========================================================================
\begin{table*}[t]
\centering
\caption{Procrustes distance per model, layer, and format pair.  Lower values indicate greater representational similarity after optimal rotation.}
\label{tab:procrustes_complete}
\small
\begin{tabular}{ll rrr r}
\toprule
\textbf{Model} & \textbf{Layer} & \textbf{SVG--TikZ} & \textbf{SVG--Asy} & \textbf{TikZ--Asy} & \textbf{Mean} \\
\midrule
Qwen2.5-Coder & L4  & 0.073 & 0.095 & 0.139 & 0.102 \\
Qwen2.5-Coder & L8  & 0.079 & 0.095 & 0.146 & 0.107 \\
Qwen2.5-Coder & L12 & 0.090 & 0.106 & 0.217 & 0.138 \\
Qwen2.5-Coder & L16 & 0.113 & 0.134 & 0.226 & 0.158 \\
Qwen2.5-Coder & L20 & 0.113 & 0.128 & 0.198 & 0.146 \\
Qwen2.5-Coder & L24 & 0.124 & 0.128 & 0.219 & 0.157 \\
Qwen2.5-Coder & L28 & 0.117 & 0.124 & 0.220 & 0.153 \\
\midrule
VisCoder2      & L4  & 0.072 & 0.094 & 0.141 & 0.102 \\
VisCoder2      & L8  & 0.080 & 0.095 & 0.150 & 0.108 \\
VisCoder2      & L12 & 0.095 & 0.108 & 0.222 & 0.142 \\
VisCoder2      & L16 & 0.115 & 0.133 & 0.222 & 0.157 \\
VisCoder2      & L20 & 0.115 & 0.133 & 0.195 & 0.148 \\
VisCoder2      & L24 & 0.123 & 0.136 & 0.216 & 0.158 \\
VisCoder2      & L28 & 0.110 & 0.125 & 0.202 & 0.146 \\
\midrule
Qwen2.5-Base   & L4  & 0.071 & 0.093 & 0.136 & 0.100 \\
Qwen2.5-Base   & L8  & 0.077 & 0.091 & 0.141 & 0.103 \\
Qwen2.5-Base   & L12 & 0.083 & 0.097 & 0.201 & 0.127 \\
Qwen2.5-Base   & L16 & 0.093 & 0.113 & 0.193 & 0.133 \\
Qwen2.5-Base   & L20 & 0.096 & 0.112 & 0.177 & 0.128 \\
Qwen2.5-Base   & L24 & 0.104 & 0.112 & 0.195 & 0.137 \\
Qwen2.5-Base   & L28 & 0.065 & 0.084 & 0.133 & 0.094 \\
\bottomrule
\end{tabular}
\end{table*}

%% ========================================================================
%% Table A3: PWCCA Values
%% ========================================================================
\begin{table*}[t]
\centering
\caption{Projection-Weighted CCA (PWCCA; \citealt{morcos2018pwcca}) mean similarity per model and layer.  Higher values indicate stronger cross-format alignment.}
\label{tab:pwcca_complete}
\small
\begin{tabular}{ll r}
\toprule
\textbf{Model} & \textbf{Layer} & \textbf{PWCCA} \\
\midrule
Qwen2.5-Coder & L4  & 0.7699 \\
Qwen2.5-Coder & L8  & 0.7734 \\
Qwen2.5-Coder & L12 & 0.8128 \\
Qwen2.5-Coder & L16 & 0.8194 \\
Qwen2.5-Coder & L20 & 0.7989 \\
Qwen2.5-Coder & L24 & 0.7786 \\
Qwen2.5-Coder & L28 & 0.7791 \\
\midrule
VisCoder2      & L4  & 0.7680 \\
VisCoder2      & L8  & 0.7743 \\
VisCoder2      & L12 & 0.8046 \\
VisCoder2      & L16 & 0.8086 \\
VisCoder2      & L20 & 0.7885 \\
VisCoder2      & L24 & 0.7764 \\
VisCoder2      & L28 & 0.7712 \\
\midrule
Qwen2.5-Base   & L4  & 0.7587 \\
Qwen2.5-Base   & L8  & 0.7657 \\
Qwen2.5-Base   & L12 & 0.8012 \\
Qwen2.5-Base   & L16 & 0.8068 \\
Qwen2.5-Base   & L20 & 0.7932 \\
Qwen2.5-Base   & L24 & 0.7760 \\
Qwen2.5-Base   & L28 & 0.8114 \\
\bottomrule
\end{tabular}
\end{table*}

%% ========================================================================
%% Table A4: Permuted-Label Probe Results
%% ========================================================================
\begin{table*}[t]
\centering
\caption{Format-classification probe accuracy with real vs.\ permuted labels.  The large gap ($\approx 0.667$) confirms that probe accuracy reflects genuine format information rather than memorized structure.}
\label{tab:permuted_probe}
\small
\begin{tabular}{ll rrr}
\toprule
\textbf{Model} & \textbf{Layer} & \textbf{Real Acc.} & \textbf{Permuted Acc.} & \textbf{Gap} \\
\midrule
Qwen2.5-Coder & L4  & 1.000 & 0.333 & 0.667 \\
Qwen2.5-Coder & L8  & 1.000 & 0.333 & 0.667 \\
Qwen2.5-Coder & L12 & 1.000 & 0.333 & 0.667 \\
Qwen2.5-Coder & L16 & 1.000 & 0.333 & 0.667 \\
Qwen2.5-Coder & L20 & 1.000 & 0.333 & 0.667 \\
Qwen2.5-Coder & L24 & 1.000 & 0.333 & 0.667 \\
Qwen2.5-Coder & L28 & 1.000 & 0.333 & 0.667 \\
\midrule
VisCoder2      & L4  & 1.000 & 0.333 & 0.667 \\
VisCoder2      & L8  & 1.000 & 0.333 & 0.667 \\
VisCoder2      & L12 & 1.000 & 0.333 & 0.667 \\
VisCoder2      & L16 & 1.000 & 0.333 & 0.667 \\
VisCoder2      & L20 & 1.000 & 0.333 & 0.667 \\
VisCoder2      & L24 & 1.000 & 0.333 & 0.667 \\
VisCoder2      & L28 & 1.000 & 0.333 & 0.667 \\
\midrule
Qwen2.5-Base   & L4  & 1.000 & 0.333 & 0.667 \\
Qwen2.5-Base   & L8  & 1.000 & 0.333 & 0.667 \\
Qwen2.5-Base   & L12 & 1.000 & 0.333 & 0.667 \\
Qwen2.5-Base   & L16 & 1.000 & 0.333 & 0.667 \\
Qwen2.5-Base   & L20 & 1.000 & 0.333 & 0.667 \\
Qwen2.5-Base   & L24 & 1.000 & 0.333 & 0.667 \\
Qwen2.5-Base   & L28 & 1.000 & 0.333 & 0.667 \\
\bottomrule
\end{tabular}
\end{table*}

%% ========================================================================
%% Table A5: A2 Permutation Test Summary
%% ========================================================================
\begin{table*}[t]
\centering
\caption{A2 permutation test ($n{=}5{,}000$): observed cross-format CKA mean vs.\ null distribution obtained by permuting sample identities within each format.  $p{<}0.001$ for all models confirms that cross-format representational similarity is not an artifact of format-level statistics.}
\label{tab:a2_permutation}
\small
\begin{tabular}{l rrrr}
\toprule
\textbf{Model} & \textbf{Obs.\ CKA Mean} & \textbf{Null Mean} & \textbf{Null Std} & \textbf{$p$-value} \\
\midrule
Qwen2.5-Coder & 0.123 & 0.041 & 0.003 & ${<}\,0.001$ \\
VisCoder2      & 0.124 & 0.044 & 0.003 & ${<}\,0.001$ \\
Qwen2.5-Base   & 0.101 & 0.037 & 0.003 & ${<}\,0.001$ \\
\bottomrule
\end{tabular}
\end{table*}

%% ========================================================================
%% Algorithm: Permutation Test Pseudocode
%% ========================================================================

\begin{algorithm}[h!]
\caption{A2 Permutation Null Test for Cross-Format CKA}
\label{alg:perm_test}
\begin{algorithmic}[1]
\REQUIRE Hidden-state matrices $\mathbf{H}^{(m,f,l)} \in \mathbb{R}^{N \times d}$ for model $m$, format $f$, layer $l$
\REQUIRE Number of permutations $B = 5{,}000$
\STATE \textbf{H$_0$:} Cross-format CKA arises solely from format-level statistics; sample identity carries no aligned signal across formats.
\FOR{each model $m$}
  \STATE Compute observed CKA for each $(l, f_1, f_2)$ pair from aligned Gram matrices
  \STATE $\bar{c}_\text{obs} \leftarrow \text{mean over all layers and pairs}$
  \FOR{$b = 1$ to $B$}
    \STATE Draw random permutation $\pi$ of $\{1, \ldots, N\}$
    \FOR{each layer $l$ and pair $(f_1, f_2)$}
      \STATE Permute rows and columns of $\mathbf{K}^{(f_2)}$ by $\pi$: $\tilde{\mathbf{K}}^{(f_2)}_{ij} = \mathbf{K}^{(f_2)}_{\pi(i),\pi(j)}$
      \STATE Compute $\text{CKA}(\mathbf{K}^{(f_1)}, \tilde{\mathbf{K}}^{(f_2)})$
    \ENDFOR
    \STATE $\bar{c}_\text{null}^{(b)} \leftarrow \text{mean over all layers and pairs}$
  \ENDFOR
  \STATE $p \leftarrow \frac{1}{B} \sum_{b=1}^{B} \mathbf{1}[\bar{c}_\text{null}^{(b)} \geq \bar{c}_\text{obs}]$
\ENDFOR
\ENSURE Per-model $p$-value, null distribution $\{\bar{c}_\text{null}^{(b)}\}_{b=1}^B$, observed/null ratio
\end{algorithmic}
\end{algorithm}

%% ========================================================================
%% Table A6: Format Probe Accuracy (moved from main text)
%% ========================================================================
\begin{table}[h!]
\centering
\small
\caption{Format classification probe accuracy (5-fold CV, logistic regression). All 21 model--layer cells are 1.0000.}
\label{tab:probe_accuracy}
\resizebox{\columnwidth}{!}{%
\begin{tabular}{lccccccc}
\toprule
Model & L4 & L8 & L12 & L16 & L20 & L24 & L28 \\
\midrule
Qwen2.5-Coder  & 1.0000 & 1.0000 & 1.0000 & 1.0000 & 1.0000 & 1.0000 & 1.0000 \\
VisCoder2       & 1.0000 & 1.0000 & 1.0000 & 1.0000 & 1.0000 & 1.0000 & 1.0000 \\
Qwen2.5-Base    & 1.0000 & 1.0000 & 1.0000 & 1.0000 & 1.0000 & 1.0000 & 1.0000 \\
\bottomrule
\end{tabular}}
\end{table}

%% ========================================================================
%% Table A7: PWCCA Null Calibration (moved from main text)
%% ========================================================================
\begin{table}[h!]
\centering
\small
\caption{PWCCA permutation null calibration (300 perms). The observed/null gap is statistically significant but small (ratio ${\sim}1.1\times$), compared to CKA's ${\sim}3\times$ null ratio.}
\label{tab:pwcca_null}
\resizebox{\columnwidth}{!}{%
\begin{tabular}{lcccc}
\toprule
Model & Observed & Null (mean$\pm$std) & Ratio & $p$ \\
\midrule
Qwen2.5-Coder  & 0.790 & 0.688$\pm$0.004 & 1.148$\times$ & $<$0.001 \\
VisCoder2       & 0.785 & 0.685$\pm$0.004 & 1.145$\times$ & $<$0.001 \\
Qwen2.5-Base    & 0.788 & 0.699$\pm$0.004 & 1.127$\times$ & $<$0.001 \\
\bottomrule
\end{tabular}}
\end{table}

%% ========================================================================
%% Table A8: SBERT Threshold Sensitivity (C3)
%% ========================================================================
\begin{table}[h!]
\centering
\small
\caption{SBERT threshold sensitivity (C3): mean CKA at L28 by matching threshold. All 252 triples have min cosine $\geq 0.70$, so thresholds $\leq 0.70$ yield identical pools. At thresholds $\geq 0.75$, $N$ drops sharply, making CKA estimates unreliable due to $n \ll d$ inflation \cite{nguyen2021wide}; we do not interpret these values as evidence of stronger alignment. The relative ordering TikZ--Asy $>$ SVG--Asy $>$ SVG--TikZ is preserved at $N{=}51$, consistent with the token-overlap gradient (F3).}
\label{tab:sbert_sensitivity}
\begin{tabular}{rr ccc}
\toprule
Threshold & $N$ & Coder & VisCoder2 & Base \\
\midrule
0.60 & 252 & 0.149 & 0.137 & 0.076 \\
0.65 & 252 & 0.149 & 0.137 & 0.076 \\
0.70 & 252 & 0.149 & 0.137 & 0.076 \\
0.75 &  51 & 0.252 & 0.247 & 0.146 \\
0.80 &   6 & 0.739 & 0.766 & 0.616 \\
\bottomrule
\end{tabular}
\end{table}

%% ========================================================================
%% Stage A Generation Details (moved from Methods for space)
%% ========================================================================
\paragraph{Stage A generation details.}
Eval pool sources: SVG from VisPlotBench and VCM-SVG; TikZ from DaTikZ v1 test; Asymptote from VisPlotBench and VCM-Asy.
All sources filtered uniformly: code length $\in [50, 8000]$ characters, caption length $\in [10, 3000]$ characters, render success within 30s, normalized-code-hash deduplication.
Generation uses vLLM \cite{vllm} in fp16 with $T{=}0.3$, $\text{top\_p}{=}1.0$, $n{=}1$, format-specific max tokens (SVG/TikZ: 1024, Asy: 2048).
3-shot exemplars are drawn from independent training data disjoint from both eval and probe pools.
The CLIPScore SNR gate was calibrated in a pilot study ($n{=}20$, 5 dimensions) using a shuffle-permutation null.

%% ========================================================================
%% Table A9: Stage A Generation Summary
%% ========================================================================
\begin{table}[h!]
\centering
\small
\caption{Stage A generation summary: validity rate per model and format (3,600 total, 0-shot and 3-shot pooled). Temperature $T{=}0.3$, $n{=}1$ per prompt. (Table~A9)}
\label{tab:stage_a_summary}
\resizebox{\columnwidth}{!}{%
\begin{tabular}{lccc}
\toprule
Model & SVG & TikZ & Asy \\
\midrule
Qwen2.5-Coder  & 1.000 & 0.995 & 1.000 \\
VisCoder2       & 0.972 & 0.962 & 1.000 \\
Qwen2.5-Base    & 0.662 & 0.722 & 0.752 \\
\midrule
Overall         & \multicolumn{3}{c}{3,227/3,600 valid (89.6\%)} \\
\bottomrule
\end{tabular}}
\end{table}

\paragraph{PWCCA in the $n \ll d$ regime.}
The limited discriminative power of PWCCA in this study reflects known properties of CCA when sample size is much smaller than feature dimensionality.
With $n{=}252$ samples and $d{=}3{,}584$ features (ratio 0.07), canonical correlations converge toward 1.0 regardless of true signal strength, compressing the dynamic range available to distinguish signal from noise.
This is the same degeneracy that motivated our removal of raw CCA from the analysis pipeline.
PWCCA corroborates the \emph{direction} of the CKA signal---observed values exceed the permutation null for all three models---but its observed/null ratio (1.13--1.15$\times$) is far below CKA's (${\sim}3\times$), limiting its value as an independent robustness check in this experimental design.

%% ========================================================================
%% C1 Negative Control: Python-X vs Visual-X Magnitude
%% ========================================================================

\subsection{C1 Negative Control: Python-X vs.\ Visual-X Magnitude}
\label{sec:c1-magnitude}

Table~\ref{tab:c1_python_magnitude} reports the per-cell comparison underlying the C1 pre-registered criterion (D060): mean python-X CKA must be $< 0.5\times$ mean visual-X CKA at each (model, layer) cell.
The criterion fails at 20 of 21 cells; only Qwen2.5-base at L28 passes, driven by the sharp final-layer washout (F4) that depresses both signals.

\begin{table}[!ht]
\centering
\small
\caption{C1 negative control magnitude comparison. Python-X = mean CKA across (python--svg, python--tikz, python--asy); Visual-X = mean CKA across (svg--tikz, svg--asy, tikz--asy). Ratio = Python-X\,/\,Visual-X. M1 criterion requires Ratio $< 0.50$; \textbf{bold} = the sole passing cell. \emph{Note:} $p$-values report the minimum across three format pairs per cell (uncorrected for multiplicity).}
\label{tab:c1_python_magnitude}
\begin{tabular}{llccccc}
\toprule
Model & Layer & Py-X CKA & Vis-X CKA & Ratio & Min $p$ (Py, uncorr.) & M1 \\
\midrule
Coder-7B & L4  & 0.093 & 0.087 & 1.06 & .015 & \ding{55} \\
Coder-7B & L8  & 0.088 & 0.091 & 0.98 & .056 & \ding{55} \\
Coder-7B & L12 & 0.092 & 0.121 & 0.76 & .033 & \ding{55} \\
Coder-7B & L16 & 0.089 & 0.141 & 0.63 & .064 & \ding{55} \\
Coder-7B & L20 & 0.084 & 0.128 & 0.65 & .198 & \ding{55} \\
Coder-7B & L24 & 0.092 & 0.141 & 0.65 & .079 & \ding{55} \\
Coder-7B & L28 & 0.102 & 0.149 & 0.69 & .005 & \ding{55} \\
\midrule
VisCoder2-7B & L4  & 0.093 & 0.088 & 1.05 & .026 & \ding{55} \\
VisCoder2-7B & L8  & 0.091 & 0.095 & 0.97 & .048 & \ding{55} \\
VisCoder2-7B & L12 & 0.099 & 0.132 & 0.75 & .047 & \ding{55} \\
VisCoder2-7B & L16 & 0.096 & 0.143 & 0.67 & .049 & \ding{55} \\
VisCoder2-7B & L20 & 0.086 & 0.130 & 0.66 & .182 & \ding{55} \\
VisCoder2-7B & L24 & 0.087 & 0.145 & 0.60 & .065 & \ding{55} \\
VisCoder2-7B & L28 & 0.091 & 0.137 & 0.66 & .023 & \ding{55} \\
\midrule
Qwen2.5-base & L4  & 0.094 & 0.085 & 1.11 & .047 & \ding{55} \\
Qwen2.5-base & L8  & 0.087 & 0.089 & 0.98 & .103 & \ding{55} \\
Qwen2.5-base & L12 & 0.088 & 0.112 & 0.78 & .110 & \ding{55} \\
Qwen2.5-base & L16 & 0.079 & 0.115 & 0.68 & .127 & \ding{55} \\
Qwen2.5-base & L20 & 0.071 & 0.109 & 0.65 & .319 & \ding{55} \\
Qwen2.5-base & L24 & 0.075 & 0.121 & 0.62 & .242 & \ding{55} \\
Qwen2.5-base & L28 & \textbf{0.030} & \textbf{0.076} & \textbf{0.39} & .236 & \textbf{\ding{51}} \\
\bottomrule
\end{tabular}
\end{table}

\paragraph{Multiplicity note.}
The 9/21 significant cells reported above use per-pair permutation $p$-values without multiplicity correction across the three format pairs tested at each (model, layer) cell.
Under Bonferroni correction for three comparisons per cell ($\alpha = 0.0167$), the Qwen family retains 2/21 significant python-X cells (Coder L4 $p{=}.015$ and L28 $p{=}.005$) versus Llama-3-8B 0/21.
The qualitative split in detectability therefore holds under correction, though the effect size shrinks from 9/21 to 2/21.

%% ========================================================================
%% Code-Naive Baseline: Llama-3-8B
%% ========================================================================

\subsection{Code-Naive Baseline: Llama-3-8B}
\label{sec:llama-baseline}

\paragraph{Note:} The initial three-layer scan was exploratory (D075); a subsequent 7-layer test with pre-registered criteria (D078: visual $\geq$70\% \& python $\leq$40\% significance) confirmed the pattern.

To test whether the cross-format CKA signal requires code-specialized pretraining, we apply the same probing framework to Meta-Llama-3-8B (base), a general-purpose LLM with no code-specific training.
We extract hidden states at 7 equidistant layers (L4/L8/L12/L16/L20/L24/L28 out of 32, corresponding to 12.5--87.5\% depth) for all 252 triples and 252 Python snippets, computing CKA with 5{,}000 permutations and 1{,}000 bootstrap subsamples.

\begin{table}[!ht]
\centering
\small
\caption{Llama-3-8B baseline CKA across 7 equidistant layers: visual cross-format pairs (top) and python--visual pairs (bottom). Bold $p$-values indicate $p < 0.005$. \emph{Note:} $p$-values are per-pair permutation tests (5{,}000 permutations, uncorrected).}
\label{tab:llama_baseline}
\begin{tabular}{lccccccc}
\toprule
Pair & L4 & L8 & L12 & L16 & L20 & L24 & L28 \\
\midrule
\multicolumn{8}{l}{\emph{Visual cross-format pairs (21/21 significant, all $p < 0.005$)}} \\
SVG--TikZ   & .053 & .071 & .078 & .090 & .092 & .089 & .088 \\
SVG--Asy    & .075 & .082 & .087 & .103 & .093 & .089 & .093 \\
TikZ--Asy   & .126 & .181 & .189 & .208 & .191 & .183 & .174 \\
\midrule
\multicolumn{8}{l}{\emph{Python--visual pairs (1/21 significant: Py--Asy L8, $p{=}.039$)}} \\
Py--SVG  & .060 & .059 & .061 & .069 & .059 & .051 & .054 \\
         & \scriptsize{.184} & \scriptsize{.517} & \scriptsize{.446} & \scriptsize{.435} & \scriptsize{.644} & \scriptsize{.652} & \scriptsize{.701} \\
Py--TikZ & .085 & .098 & .103 & .102 & .094 & .084 & .087 \\
         & \scriptsize{.182} & \scriptsize{.151} & \scriptsize{.160} & \scriptsize{.189} & \scriptsize{.130} & \scriptsize{.119} & \scriptsize{.096} \\
Py--Asy  & .107 & .115 & .116 & .118 & .092 & .083 & .093 \\
         & \scriptsize{.118} & \scriptsize{\textbf{.039}} & \scriptsize{.146} & \scriptsize{.221} & \scriptsize{.310} & \scriptsize{.218} & \scriptsize{.127} \\
\bottomrule
\end{tabular}
\end{table}

Visual cross-format CKA is significant at all 21 cells ($p \leq 0.0024$; mean 0.116).
The TikZ--Asy $>$ SVG--Asy $>$ SVG--TikZ ordering replicates at all 7 layers, strengthening the token-overlap gradient (F3) as robust to the absence of code-specialized pretraining.
In contrast, python--visual pairs are overwhelmingly non-significant (20/21 cells $p > 0.05$); the sole exception is Python--Asy at L8 ($p = 0.039$), a borderline result that does not survive multiplicity correction.

We note that the equidistant protocol samples Llama-3 to L28 (87.5\% depth), missing the final layer (L31, 97\% depth) where D075's three-layer scan found the only two significant python cells.
Thus, python--visual entanglement in code-naive Llama-3, if present, is \emph{layer-local} (present only at the deepest tested layer in D075, L31, 97\% depth), whereas in code-pretrained Qwen models, python-X significance is broader (9/21 cells $p<0.05$, uncorrected) with observed-to-null ratios of ${\sim}1.3$--$1.7\times$ versus Llama's $0.95$--$1.09\times$.
Importantly, the magnitude pattern is shared: Llama's python/visual CKA ratios (0.60--0.99) overlap with the Qwen range (0.60--1.11).
The split is thus in \emph{permutation detectability}---code pretraining elevates python-X CKA above the permutation null---not in raw magnitude.
